\documentclass[../main.tex]{subfiles}
\begin{document}
%==========================================TIÊU ĐỀ TẬP========================================
\begin{table}[h]
  \begin{adjustwidth}{-1cm}{}
    \begin{tabular}{>{\centering\arraybackslash}p{6cm}|>{\raggedright\arraybackslash}p{12.5cm}}
      \multirow{5}{6cm}
      % Cột bên trái: ÉPISODE và số tập
      {\\[-40pt]
        \flaregothic\fontsize{20pt}{20pt}\selectfont \centering
        {\contour{errorfive}{\textcolor{black}{ÉPISODE}}}\color{black}\\
        \burbank\fontsize{72pt}{72pt}\selectfont
        \includegraphics[height=3.12359cm]{../../../../Image/Book?/number/num27.png}
        \\[18pt]
        % \flaregothic\fontsize{14pt}{14pt}\selectfont
        % \color{epattr}BIENVENUE, \uline{\color{Banpire} Mel} !
        % \flaregothic\fontsize{18pt}{18pt}\selectfont
        % \color{epattr}ÉPISODE PERDU
        \color{black}
      }
       &
      % Dòng thứ nhất: tên tập tiếng Nhật
      {\fotskip\fontsize{18pt}{18pt}\selectfont \ruby{新}{あたら}\ruby{入}{にゅう}\ruby{生}{せい}の\ruby{歓}{かん}\ruby{迎}{げい}\ruby{会}{かい}}
      \\[6pt]
      % Dòng thứ hai: tên tập tiếng Anh
       & {\cafeteria\fontsize{18pt}{18pt}\selectfont New Students' Welcome}                                                                                                                          \\[4pt]
      % Dòng thứ ba: tên tập tiếng Việt
       & {\vnmsans\fontsize{18pt}{18pt}\selectfont Bữa tiệc chào đón học sinh mới}                                                                                                                   \\[8pt]
      % Dòng thứ tư: Nhân vật xuất hiện
       & {\caviar{\fontsize{14pt}{14pt}\selectfont Nhân vật: \emp{choco} \emp{azki} \emp{haato} \emp{aqua} \emp{shion} \emp{ayame} \emp{watame} \emp{matsuri} \emp{fubuki} \emp{kanata} \emp{towa}}}
      % \\[6pt]
      % Dòng cuối cùng: Thời gian diễn ra của tập (thường sẽ là ngày phát hành)
      % & {\continuum\fontsize{16pt}{16pt}\selectfont Ngày phát hành: 11/06/2022}\\[8pt]
    \end{tabular}
  \end{adjustwidth}
\end{table}
%===========================================NỘI DUNG==========================================
\color{black}

\begin{center}
  \uline{Hai ngày sau.}
  \pov{Hikawa Kaigou}
\end{center}

Hai ngày trôi qua kể từ tiết Sinh học nhộn nhịp ngày hôm đó. Tôi đã bắt đầu quen dần với nhịp học ở lớp 1-C.

Không hiểu sao, cái lớp này hợp với tôi hơn tôi tưởng. Có lẽ bởi vì trong lớp không có một mống đàn ông nào, ngoài Kuon-kun ra, nên có lẽ tôi cảm thấy thoải mái hơn?

Con gái trong lớp dễ bắt chuyện, mà tính tôi thì vốn cũng dễ kéo người khác lại gần, miễn là đối phương tốt với tôi. Chỉ cần vài câu hỏi xã giao, vài lần đổi chỗ làm bài nhóm, thế là tôi đã kết thân được một nửa lớp rồi. Nhất là mấy bạn năng lượng cao như Hanabi-san, Kana-san, hay Hotaru-san. Họ nói chuyện với họ y như trượt vào dòng nước chảy vậy, không cần cố gắng quá mức.

Chỉ trừ một vài bạn trong lớp thì\dots\ hơi kỳ quặc một chút, như Nanase-san cứ xa lánh tôi - tôi cá là bởi vì cú tát trời giáng đó, mà tôi cũng không nghĩ là cô ấy còn nhớ. Hay Honoka-san và Suzune-san, những người tự biến mình thành một gánh xiếc của lớp.

Với tôi là như vậy, còn Suzu thì khác. Nó nhẹ nhàng, lễ phép, nói câu nào ra câu nấy; nhờ thế mà mấy bạn trong lớp rất nhanh quý nó. Mấy bạn nữ còn kiểu ``em gái nhỏ đáng yêu'', lúc nào cũng miệng tươi cười mỗi khi thấy con bé đi ngang. Con bé thì khụng khiệng như chú mèo con ngoan hiền, lúc nào cũng xin lỗi xin phép, đúng kiểu khiến người ta muốn bảo vệ.

Kuon-kun thì\dots\ Hừm. Nói sao nhỉ.

Cậu ta kiểu như ``thích ứng nhanh đến mức đáng sợ''. Vào lớp chưa đầy tuần nhưng đã sống như thể đây là trường của mình từ hai năm trước. Bạn bè thì cậu không chủ động tìm, nhưng ai hỏi gì thì vẫn trả lời đàng hoàng. Mà do mặt cậu ta lúc bình thường trông như chẳng quan tâm đến chuyện gì, nhiều đứa lại nghĩ cậu ``ngầu cool'' kiểu bí ẩn.

Nên rốt cuộc, bằng cái thái độ nửa hờ hững nửa lịch sự của mình, cộng với việc là con trai duy nhất của lớp, Kuon-kun lại được khá nhiều người để ý.

Thú thật, đôi khi tôi cũng thấy chạnh lòng. Mình phải tốn ba câu chuyện mới kết được một người bạn, còn Kuon-kun chỉ cần đứng yên ba giây đã có người tới làm quen.

Còn về cuộc trò chuyện giữa Kuon và Yuko-sensei hai ngày trước đó\dots

Tôi với Suzu không nghe lén. Lại càng chẳng có khả năng nghe xuyên tường như mấy nhân vật trong manga.

Tối hôm ấy, khi ba chúng tôi đang thu dọn đống sách vở lỉnh kỉnh sau bữa tối, tôi cố tình hỏi vu vơ: ``Kuon-kun, Yuko-sensei giữ cậu lại nói gì thế?''

Cậu ta dừng tay lau nhà, im lặng đúng một hơi thở, rồi đáp như thể chỉ kể chuyện trời nắng trời mưa: ``Cô ấy hỏi về câu lạc bộ\dots\ và về `chuyện ngoài sân' hôm đó.''

Cây bút máy của Suzu lập tức rơi xuống sàn. Tôi thì cứng đờ, tim đập trật nhịp.

Chúng tôi chưa bao giờ hẹn ước, nhưng cả ba đều nhớ như in:

Hình ảnh bóng đứa bé kimono xanh thẫm lủi thủi giữa sân đá.

Bầu không khí trong lớp bỗng dưng lạnh toát, đủ để hơi thở thành sương.

Cùng với đó là những tiếng khóc trẻ con thoảng qua, mỏng như lá kẽm cào vào ống tai.

Và, cảm giác có thứ gì đang giật lấy mảnh ký ức chưa kịp vá của chúng tôi.

Tôi nuốt khan, giọng khàn đặc: ``Sensei\dots\ biết được bao nhiêu?''

Cậu ta khẽ nhún vai, ánh mắt xa xăm: ``Đủ để không gọi là `nhìn nhầm'.''

Một luồng điện râm ran chạy dọc sống lưng tôi. Không phải nỗi sợ ma quỷ, mà là cái cảm giác bất lực khi nhận ra: cánh cổng Aogami vừa khép lại phía sau chúng tôi đã biến thành bức tường không lối về.

Mọi thứ\dots\ dường như không còn đơn giản như ngày đầu mới tới Aogami nữa.

Nhưng, trong lòng tôi giờ lại có một nỗi tâm tư khác.

Tôi vẫn nhớ như in khoảnh khắc trong The Void - nơi thời gian bị bóp méo, xung quanh chúng tôi chỉ có một màu đen bóng tối - khi Kuon-kun thuyết phục hai chị em tôi sống ở thế giới của cậu.

Lúc đó tôi chỉ nghĩ đơn giản:

``Thoát khỏi cái vòng lặp chết tiệt.''

``Sống lại.''

``Tìm một điểm bắt đầu mới.''

Nhưng bây giờ, khi ở Aogami thực, một thế giới có lịch sử, có nhịp sống, có người, có bạn bè\dots\ Tôi lại không chắc nữa. Đặc biệt là sau những gì đã xảy ra với chúng tôi ở thế giới Cầu Vồng trước và những vòng lặp tưởng chừng vô tận đó.

Đây có thật sự là nơi quyết định để bắt đầu lại không? Hay chỉ là một lớp ký ức khác của cái vòng xoáy mà chúng tôi chưa hiểu nổi?

Câu trả lời vẫn mơ hồ.

Cứ như thể tôi đang đứng một chân trên mép vực, chân còn lại chưa biết điểm đặt.

Trong lúc đầu tôi vẫn đang chìm nghỉm với những câu hỏi tu từ không có lời giải đó, bất chợt\dots

``\textbf{Kaigou-chaaaaan! Bọn tôi có việc này muốn nhờ cậu được không\~{}?}'' Tiếng Hanabi-san bỗng vang lên bên tai tôi như bóp nổ.

Tôi ngẩng đầu.

Cô nàng cùng Kana-san đã dán sát mặt xuống bàn ba đứa tôi, mắt lấp lánh như thể sắp mở màn kể chuyện trên trời dưới biển.

Đằng sau, Inari-san mỉm cười lịch sự, còn Kanade-san ôm khư khư xấp bài như đang giữ cả sinh mệnh học tập.

Tôi kịp chớp mắt, rồi cười: ``Ừ, kéo ghế ngồi đi. Đúng hôm cần người tám chuyện.'' Tạm thời gác lại mớ suy tư vừa nãy.

Liếc đồng hồ treo tường, còn chừng năm phút nữa mới hết giờ ra chơi.

Cô bạn tóc nâu vừa hạ mình xuống ghế, Kana đã lượn sang đối diện, cằm gác tay, nhìn chúng tôi như thám tử chuẩn bị lấy lời khai.

``Được rồi, bọn tớ xung phong trước nhé!'' Hanabi reo lên, mắt long lanh như đèn sân khấu.

Inari giật mình, khẽ giật tay áo bạn: ``Ha-Hanabi-chan\dots\ đừng dọa họ, hỏi nhẹ nhàng thôi.''

Kanade thở dài, ngồi sát cô ấy như người giữ phanh, ngăn hai cô kia lao xuống dốc không phanh.

Tôi cười nhẹ: ``Không sao đâu. Cứ hỏi đi, tôi cũng tò mò các cậu muốn gợi chuyện gì.''

Suzu gật nhẹ: ``Bọn mình không phiền.''

Kuon-kun dựa lưng, gật đầu: ``Cứ hỏi, đừng dông dài là được. Mà các cậu hỏi bọn tôi có việc gì vậy?''

Hanabi vỗ tay, đuôi ngựa khẽ lắc: ``Thật ra bọn tớ đang lên kế hoạch tổ chức tiệc chào mừng học sinh mới, nên phải khảo sát trước khẩu vị, sở thích, món ưa thích của mọi người để đặt thực đơn cho chuẩn!''

``Ra thế. Rồi, hỏi gì, hỏi luôn đi,'' tôi nói, khoanh tay hất cằm về bốn cô bạn, ra vẻ hối thúc.

\il

Cô bạn bờm cáo trắng bắt đầu nhỏ nhẹ: ``Vậy\dots\ món ăn yêu thích của mọi người là gì?''

Tôi trả lời trước: ``Cái gì miễn bỏ bụng không chết là được.''

Kuon-kun thì xen vào: ``Tôi chọn thịt bò. Miễn là thịt nói chung là ổn. Không thích ăn rau lắm, ngoại trừ một vài loại.''

Còn Suzu nghĩ một lát rồi liền đáp: ``Mình thích món thanh đạm. Cá hay trứng chẳng hạn.''

``Khác hẳn ba người kia,'' Kana-san cười. ``Dễ thương quá.'' Nghe vậy, con bé liền đỏ mặt cúi xuống.

Tôi nghiêng đầu cười, đúng kiểu khiến người ta muốn trêu thêm.

Kanade-san ghi bút: ``Còn môn thể thao ưa thích?''

Kuon-kun chống cằm, ngẫm lâu rồi nói: ``\dots Chắc chạy bộ.''

Tôi huých khuỷu tay: ``Cậu mà chạy nổi? Lười chảy thây thế kia.''

Cậu ta đáp lại, giọng thờ ơ: ``Tôi chạy được, chỉ không thích thôi.''

Kana-san cười phá lên. ``Cậu thì sao, Kaigou-chan?''

Tôi nhún vai: ``Tôi chơi được hết. Đừng hỏi môn ghét, còn môn thích thì\dots\ nhiều lắm.''

Inari-san nhìn tôi, mắt hơi ngưỡng mộ. ``Còn Suzuki-san thì sao?''

Em tôi lí nhí: ``Em mê bóng bàn\dots\ vì không phải chạy nhiều\dots ''

``Lý do thật luôn hả?'' tôi hỏi, nó gật nhẹ, cả bàn bật cười.

``Này này! Câu đinh!'' Hanabi-san nghiêng người. ``Vậy thì\dots\ Game yêu thích?''

Tôi và Kuon đồng thanh:``Nhạc và giải đố.''

Cô bạn bờm cáo tròn mắt: ``Ôi, hai người lại cùng có chung sở thích này!''

Suzuki cười hiền: ``Em thích game nhẹ, có chút đáng yêu\dots'' Thấy vậy, cả bốn cô gái chắp tay, cùng nhau nhìn gương mặt Suzu càng ửng lên. Mà công nhận, con bé khi xấu hổ trông dễ thương thật.

Cô bạn tóc nâu bất ngờ nghiêng đầu, mắt tò mò lạ thường. ``Vậy\dots\ điểm thu hút nhất trên người mọi người là gì?''

Cả ba chúng tôi cùng cứng đờ với câu hỏi như rơi từ trên trời xuống.

``\dots Câu này liên quan gì đến tiệc?'' tôi hỏi thẳng.

``Ơ thì\~{}'' cô ấy toe toét. ``Biết điểm mạnh nhau dễ thân hơn mà!''

Inari giật tay cô bạn, khẽ lắc đầu: ``Ha-Hanabi-chan\dots\ không hợp chút nào\dots ''

Cậu con trai thì khoanh tay, nhíu mày: ``Hơi riêng tư đấy.''

Suzu luống cuống, ậm ừ không biết trả lời thế nào.

Nhưng Hanabi-san không dừng. Cô ấy chống cằm, nhìn tôi rồi liếc sang em tôi bằng một cái nhìn soi mói, và buông ra một câu hỏi lấp lửng: ``Thế\dots\ hai cậu hôm nay mặc màu gì?''

Tôi chớp mắt: ``\dots Gì cơ?''

Cô em tôi thì lùi sát lưng ghế:``T-Tại sao\dots\ lại hỏi bọn mình như vậy? M-Mà\dots\ cậu đang hỏi tới cái gì thế?''

``Cậu biết đó\~{}'' cô ấy nhoẻn miệng cười toe, tay chỉ vào phần ngực và phần eo của mình, ``cái `nội thất' bên trong á\~{}''

Tôi hiểu ngay, rõ ràng là đang nói về cái gì. Ngay cả Kuon-kun cũng quay sang, ánh mắt tỏ rõ bất lực, miệng lầm bầm ``không ngờ rằng cô ấy ở thế giới này cũng dâm đãng như vậy.''.

Tôi nở một nụ cười, nhưng nắm tay phải siết chặt, chuẩn bị giơ lên: ``Hanabi-san. Muốn tôi đập chết cậu bây giờ không?''

``\textbf{Inari-san!}'' cô bạn Thiên sứ hét.

Cô bạn Cáo tuyết lập tức kéo cô ấy lại: ``Đ-Đừng! Không hỏi nữa! Cậu làm Kaigou-chan khó chịu đấy!'' Trong khi đó, cô bạn chỉ khúc khích, như không làm gì sai.

Rồi bất chợt, Kuon-kun lên tiếng. Giọng cậu thản nhiên\dots\ đến nguy hiểm: ``Còn các cậu? Hôm nay\dots\ màu của các cậu là gì?''

Cô bạn tóc đuôi ngựa liền giật mình, còn Inari-san đỏ bừng như bị bắt quả tang.

Kana-san vỗ tay: ``Kuon-kun phản đòn kìa!!'' còn Kanade-san đập bàn, mặt đỏ như xôi gấc: ``\textbf{Thôi!} Hai bên thôi đi! Không ai được hỏi những câu dạng này nữa!''

Tôi khoanh tay, cười đắc thắng. ``Đó. Đến các cậu còn chẳng dám trả lời huống gì tới chúng tôi.''

Suzu nhỏ giọng, gương mặt vẫn còn đỏ hỏn: ``Onii-sama\dots\ đừng hỏi vậy\dots\ em ngại\dots''

Cậu ta chỉ khoát tay: ``Chẳng ai ép em phải nói cả. Anh cũng mong đừng ai ngô nghê đến mức đó.''

Sau tràng cười, không khí bớt căng, cả bàn thở phào nhẹ nhõm; tiếng cười vẫn còn vương trên môi, nhưng ánh mắt bốn cô gái đã chuyển sang một thứ ấm áp hơn, như muốn nói lời xin lỗi bằng sự dịu dàng bất ngờ.

Inari-san thở dài, giọng trầm xuống, chân thành hẳn: ``Thật ra\dots\ bọn mình muốn hiểu ba người nhiều hơn, không chỉ vì cái tiệc thôi đâu.''

Hanabi-san cũng nghiêm mặt, lần đầu không cười đùa: ``Ừ. Để bữa tiệc vui thật sự, phải biết mọi người thích gì, ghét gì chứ.''

Kana-san thì gãi má, cười hiền: ``Hơn nữa, ba cậu mới đến\dots\ bọn tớ muốn các cậu cảm thấy đây là nhà.''

Tôi chợt khựng một nhịp. Câu nói ấy không hề xã giao. Nó không đơn thuần là lời khách sáo, mà như thể cô ấy đang mở một cánh cửa mà tôi từng nghĩ mình chỉ có thể đứng ngoài nhìn vào.

Suzu khẽ cúi đầu, má ửng hồng: ``Mình\dots\ cảm ơn mọi người.''

Kuon-kun gật nhẹ, đôi mắt vẫn lạnh nhưng giọng mềm hơn: ``Quả thật, biết trước tính nhau cũng tiện.''

Cô bạn tóc đuôi ngựa chống nạnh, tinh nghịch trở lại: ``Vậy là xong rồi! Tuần sau, bữa tiệc sẽ là phiên bản `có đặt làm theo' đấy!''

Kanade-san đập nhẹ vào vai cô bạn: ``Nhưng đừng có đặt món quái thai nào đấy nhé.''

``Biết rồi\~{}!''

Cô bạn bờm cáo cúi đầu, nhỏ nhẹ: ``Cảm ơn các cậu đã chia sẻ\dots\ và rất vui được làm quen.''

Tôi mỉm cười, lắc nhẹ cổ: ``Không có chi. Cảm ơn vì sự hiếu khách\dots\ dù hơi quá nhiệt.''

Cô bạn tóc tím cười khoái trá: ``Nhiệt thế mới vui chứ! Giờ bọn tôi hỏi những người khác đây!''

\il

Khi bốn người họ quay đi để tiếp tục phỏng vấn các bạn khác, tôi tựa lưng vào ghế, để cho ánh mắt mình trôi lơ đãng.

Không hiểu vì sao\dots\ ký ức cũ lại trồi lên.

Ở những vòng lặp trước, khi bọn tôi còn mang những cái tên khác, khi Kuon-kun còn là Kyuen, còn tôi là Saikyou, bữa tiệc chào mừng này chưa từng thật sự tồn tại.

Có những vòng lặp, nó chưa kịp diễn ra thì mọi thứ đã sụp đổ.

Có lúc, trường học bị phá hủy bởi một tai nạn không thể giải thích.

Có lúc\dots\ chúng tôi chỉ có thể chôn chân nhìn các bạn học chết mà chẳng thể làm gì được.

Có vòng, lời nguyền của Shino-san bùng phát trước cả khi ai đó kịp nhận ra điều bất thường.

Và rồi, như một vết xước bị xóa sạch, tất cả lại quay về điểm bắt đầu. Quay về ngày đầu tiên mà chúng tôi nhập học.

Tôi rùng mình khi nhớ lại cảm giác lúc đó.

Có một chu kỳ\dots\ tôi từng nghĩ Kuon đã ra đi mãi mãi. Không đơn thuần là xa cách, mà là ranh giới mong manh giữa sự sống và cái chết.

Người cậu ấy lúc ấy bê bết máu, hơi thở mong manh đến nỗi tôi phải áp tai sát ngực mới cảm nhận được. Dù cậu vẫn cố nói ``không sao'', tôi vẫn bướng bỉnh bế cậu chạy thẳng tới phòng y tế.

Ký ức ấy vẫn thỉnh thoảng bất ngờ ập đến, bóp nghẹt trái tim tôi. Chỉ vài giây ngắn ngủi thôi, nhưng đủ khiến tôi rùng mình.

Tôi từng căm ghét đàn ông, ghét đến mức không buồn đặt niềm tin vào bất kỳ ai. Thậm chí còn cấm cản việc Suzu giao du với bất cứ con trai nào, ắt cũng vì an toàn của chính con bé.

Nhưng rồi Kuon-kun xuất hiện, thay đổi mọi thứ. Cậu ta có vẻ ``bố đời'', thô kệch, thẳng thừng đến mức thiếu tế nhị, nhưng lại là người tốt tính. Cậu ta không đòi hỏi gì ở hai chị em tôi; trái lại, cậu ném cho chúng tôi một phao cứu sinh, hứa hẹn một thế giới tươi sáng hơn, dù bấy giờ vẫn chưa thể chạm tới.

Nhưng lần này, mọi thứ đã khác.

Tôi liếc nhẹ quanh lớp. Những khuôn mặt cười đùa, tranh luận, hỏi han nhau bình thường như bao học sinh khác.

Không dấu hiệu sụp đổ.

Không cảm giác ``đếm ngược''.

Không làn lạnh quen thuộc báo hiệu vòng lặp sắp khép lại.

Tôi thở ra nhẹ nhõm, khóe môi khẽ nhếch.

``Thế giới thật này\dots\ giờ đã rẽ sang trang mới.''

Chỉ một dòng nghĩ thầm, nhưng đủ để tôi tin: lần này, chúng tôi đang bước trên con đường khác biệt.

Tôi ngẩng đầu lên đúng lúc thấy nhóm bốn người ban nãy đã chuyển sang các bàn kế bên, nơi sáu học sinh mới ngồi đó.

Azuki-san là người trả lời đầu tiên.
Cô ấy nói chuyện thẳng thắn, không vòng vo, thỉnh thoảng còn chen vài câu đùa khiến Hanabi-san cười nghiêng ngả. Kiểu người này chắc chắn sẽ là trung tâm của các cuộc trò chuyện.

Koishin-san thì khoanh tay, quay mặt đi. ``Tôi không quan tâm mấy thứ này\dots\ nhưng nếu cần thì trả lời cũng được.'' Giọng lạnh tanh, nhưng ánh mắt lại liếc sang xem phản ứng người khác. Dường như mọi người không phản ứng gì quá tiêu cực với phản ứng của cô ấy, ngoại trừ những tiếng cười trừ. Quả đúng là một tsundere chính hiệu.

Aku-san thì trái ngược hoàn toàn.
Mỗi câu trả lời đều run run, tay đan vào nhau, giọng nhỏ đến mức Inari-dsn phải nghiêng sát mới nghe được. Nhưng tôi thấy ở cô ấy một sự chân thành rất rõ.

Shion-san thì\dots\ thôi, khỏi nói. ``Sở thích của ta là nghiên cứu ma thuật cấm và giải phóng sức mạnh bị phong ấn!'' Cách nói khoa trương đến mức Kana suýt sặc cười, còn Kanade phải ôm trán. Chuunibyou, nhưng là kiểu không giấu giếm.

Ayame-san thì cười từ đầu tới cuối.
``Sở thích của tôi á? Chơi khăm!'' Nói xong còn cười nhây, khiến cả bàn phải đề phòng ngay lập tức. Loại này, hoặc là gây rắc rối, hoặc là mang lại không khí cực kỳ sôi động. Chẳng thể nào đoán được.

Và cuối cùng là Wata-san. Giọng nhẹ như mây: ``Món mình thích là bánh kem\dots\ nên nếu mọi người muốn, mình có thể nướng cho cả lớp.''

Cả nhóm im lặng đúng một nhịp, rồi Hanabi-san reo lên đầy phấn khích, như thể sắp có đồ ngon.

Tôi nhìn cảnh đó, trong lòng dâng lên một cảm giác khó gọi tên. Những con người này\dots\ tính cách khác nhau, phản ứng khác nhau, chẳng ai giống ai.

Nhưng chính vì thế, tôi lại thấy có gì đó rất hứa hẹn.

Nếu thế giới này thật sự không lặp lại nữa\dots

Nếu bữa tiệc chào mừng này thật sự diễn ra\dots

\dots thì có lẽ, lớp 1-C sẽ không chỉ là một điểm dừng tạm thời.

Tôi chống cằm, tiếp tục quan sát, để cho cảm giác đó lắng lại. ``\textit{Có lẽ lần này, mọi thứ sẽ đi xa hơn những gì chúng tôi từng biết.}''

Giờ nghỉ sắp hết, tiếng ồn trong lớp trở lại. Không hỗn loạn như kiểu vòng lặp bị đứt gãy, mà là thứ âm thanh đời thường, tươi sáng, như lớp học đang thở thật.

Tôi nghiêng sang Kuon-kun và Suzu, giữ giọng ở mức chỉ ba đứa nghe được. ``\vie{Hai người thấy sao? Mấy ngày qua có gì lạ không?}''

Kuon khẽ nhún vai, ngón tay vẫn quay bút như vô thức. ``\vie{Quá yên ả, đến mức nghi ngờ. Đi học, tán gẫu, về nhà. Không bị giật lại, không có cảm giác như bị quay về từ đầu.}''

Suzuki gật nhẹ, hai tay đặt khít trên mặt bàn.
``\vie{Em cũng vậy\dots\ nụ cười của mọi người giờ mới thật sự tự nhiên. Không còn kiểu cười gượng như trước.}''

Tôi khẽ ``ừ'' một tiếng.

Đúng thế.

Từ ngày nhập học đến nay, dù chỉ vài hôm, nhịp sống đã khác. Không còn cảm giác ``hôm nay có thể là ngày cuối'', không lo sáng mai tỉnh dậy mọi thứ trở về con số không.

``\vie{Bữa tiệc chào mừng này,}'' tôi chậm giọng, ``\vie{có vẻ sẽ thành hiện thực.}''

Cậu ta liếc qua. ``\vie{Nghe như cô đang tự trấn an mình đấy nhỉ.}''

``\vie{Không hẳn,}'' tôi mỉm cười, ``\vie{chỉ là kiểm định lại thôi.}''

Suzu mỉm cười, nhưng đáy mắt em lấp lánh vẻ lo lắng: ``\vie{Em mong mọi thứ sẽ vui vẻ. Nhưng\dots\ vẫn cảm thấy có điều gì đó không ổn.}''

Tôi nhìn em, nhướng mày. ``\vie{Không ổn chỗ nào?}``

``\vie{\dots Như cảm giác trước cơn bão,}'' em thì thầm. ``\vie{Trời vẫn trong, nhưng gió đã xoay chiều.}''

Kuon lặng thinh vài nhịp.

``\vie{Dự báo à?}'' cậu ta lẩm nhẩm. ``\vie{Nghe quen quá.}''

Tôi khẽ cười ``\vie{Dù sao đi nữa, nếu chuyện gì xảy ra--}''

``\vie{Tôi biết,}'' Kuon-kun ngắt lời, giọng bình thản. ``\vie{Lần này, dù điều gì xảy ra, chúng ta sẽ đối mặt với nó. Dứt điểm một lần và mãi mãi.}''

Lời ấy khiến tôi chết lặng. Không phải bất ngờ, mà vì nó cứng rắn đến mức không chỗ cho hoài nghi. Nếu cậu ta đã quyết đinh như thế, thì tôi cũng sẽ tin, rằng lần này chúng ta sẽ không còn phải quay lại từ đầu nữa.''

\il

Tôi nghiêng người, hạ giọng xuống mức chỉ mình cậu nghe được.

``\vie{À mà này.}''

Kuon-kun liếc sang. ``\vie{Gì?}''

Tôi ghé sát tai cậu, nói rất khẽ: ``\vie{Hôm nay tôi mặc màu tím đó.}''

\dots

\dots

``\vie{\dots Hả?}''

Cậu ta đông cứng. Toàn thân cậu khựng lại đúng một giây, rồi mồ hôi hột bắt đầu lăn xuống thái dương. Gương mặt đỏ hỏn lên.

``\vie{\dots Tại sao cô lại nói cho tôi biết cái đó chứ?}'' giọng cậu ấy tỏ rõ đầy bất lực.

Tôi nheo mắt cười. ``\vie{Dựa trên cái mặt cậu hồi nãy khi Hanabi-san hỏi, tôi dám chắc là cậu cũng muốn biết đấy thôi\~{}}''

``\vie{Onee-sama\dots}'' Suzu kéo nhẹ tay áo tôi, nhắc nhở. ``\vie{Đừng trêu anh ấy chứ\dots}''

Lúc này, chiếc đồng hồ đeo tay trên cổ tay của cậu con trai khẽ rung nhẹ, một giọng nam trầm khẽ cười vang lên nhỏ trong tai tôi, chỉ mình tôi nghe được: ``\eng{Chà, chà. Không ngờ được cô cũng có ngày trêu chọc người khác thế này. Thật khó tin là cũng chính cô là người đã tè dầm trong thế giới của tôi đấy. Bây giờ lại thành cô nàng dũng cảm trêu chọc bạn trai rồi hả?}''

Bị giọng của Game làm cho giật mình, tôi suýt nữa thì sặc nước bọt. Vừa thở dài vừa mắng thầm vào micro nhỏ trong đồng hồ: ``\eng{Đã bảo là đừng nhắc tới chuyện đó rồi mà!}''

Kuon-kun thì không biết nên phản ứng thế nào, chỉ có thể quay mặt đi, lẩm bẩm cái gì đó nghe như lời cầu nguyện.

Tôi khoanh tay, tiếp tục thì thầm với cậu ấy: ``\vie{Cậu nên cảm thấy may vì đó là cậu đấy. Chứ là thằng con trai nào khác thì tôi đã cho một cái đấm ngay từ lúc được hỏi rồi.}''

Suzu, dù không nghe được tiếng Game, cũng để ý thấy tôi vừa rồi lẩm bẩm gì đó với cổ tay, bèn nghiêng đầu hỏi nhỏ: ``\vie{Game-san nói gì thế ạ??}''

``\vie{Không có gì, đừng hỏi. Có một ông cụ non thích móc mỉa người khác thôi,}'' tôi đáp, lướt mắt qua chiếc đồng hồ như muốn thôi miên nó im lặng.

Kuon-kun ngó sang, cũng thắc mắc: ``\vie{Game lại nói gì nữa à?}'' Cậu ta quen rồi với việc chiếc đồng hồ thông minh kia thỉnh thoảng xen vào câu chuyện, dù chỉ có Kaigou nghe được giọng của Game.

``\vie{Không có gì.}'' tôi bực bội thì thầm.

Suzu thở dài khe khẽ, nhưng tôi thấy khóe môi con bé cong lên một chút.

\il

Chuông vào tiết vang lên, cắt ngang cuộc nói chuyện.

Tôi ngồi thẳng lại, ánh mắt lướt qua lớp học lần nữa.

Bữa tiệc chào mừng sắp tới. Những con người mới. Một thế giới không còn lặp lại, ít nhất là chưa.

Tôi không biết điều gì đang chờ phía trước.

Nhưng có một điều tôi chắc chắn.

Lần này, nếu có chuyện xảy ra\dots\ chúng tôi sẽ đối mặt với nó trong tư thế của những người đang sống thật.

\ct

Buổi chiều hôm ấy, tôi cùng Shino-san và Sakura-san bước trên con đường tắt quen thuộc, nơi những tán cây rợp bóng và hàng rào gỗ mốc xỉnh tạo thành một hành lang yên tĩnh giữa khu dân cư.

Kuon-kun và Suzu đã đi trước; cậu ta bảo phải ghé tiệm tiện lợi nhận ca làm thêm sắp tới, còn em tôi thì dính sát như cái bóng, sợ lạc đường giữa thị trấn còn lạ lẫm. Ba người chúng tôi rẽ vào con hẻm nhỏ, nơi lá cây thường xào xạc theo gió, nhưng hôm nay, không khí tĩnh lặng đến bất thường.

Không phải mối tĩnh mịch dễ chịu, mà là thứ yên lặng khiến lỗ tai bỗng ù đặc, như thiếu một nhịp thở của trời đất.

``Gió hôm nay yếu thật,'' tôi cất tiếng, chỉ để xua đi cái cảm giác bí bách đang bóp nhẹ lấy lồng ngực.

Sakura-san khẽ ``ừ'', mắt vẫn dán vào tán cây đang đứng im như tượng. ``Đúng giờ này thường lá rừng phải động rồi chứ nhỉ.''

Shino-san không nói gì. Cô ấy bước chậm rãi, ánh nhìn dán vào bờ rào gỗ xám xịt bên đường. Phía sau đó là khu vườn hoang cỏ lút đầu gối, nơi đồn thổi từng có một ngôi nhà gỗ sụp đổ từ thời bao cấp.

Tôi vừa định lên tiếng thì\dots

*CẠCH*

Một tiếng khẽ, như có mảnh gỗ vừa chạm vào nhau trong bóng tối.

Phản xạ khiến tôi đứng chôn chân. ``\hl{Nghe thấy chứ?}'' tôi thì thầm.

Cô nàng tóc hồng gật, tay đã nắm chặt dây cặp. ``Rõ mồn một.''

Cô  nàng mắt nâu chỉ khẽ gật, mắt không rời hàng rào.

Ba chúng tôi đứng im, hơi thở cũng nhỏ như sợ làm vỡ không gian.

Không bóng người.
Không linh hồn lẩn quất.
Không một làn gió.

Chỉ có chiếc chuông gió bằng đồng xanh đang lắc lư dưới mái hiên nghiêng ngả của ngôi nhà hoang. Nó rung nhẹ, phát ra tiếng *LENG KENG* mong manh, dù lá liễu bên cạnh không hề động đậy.

Tôi thấy gáy mình tê dại.

``Cái chuông ấy\dots'' tôi lẩm nhẩm, ``trước đây tôi chưa từng thấy nó.''

Sakura-san nheo mắt, giọng trầm xuống. ``Không phải mới. Nhìn lớp bụi bám kia có lẽ đã treo đó cả thập kỷ rồi.''

Shino-san đặt tay lên ngực, tiếng nói mỏng như khói: ``Nhưng\dots\ tại sao nó lại kêu chứ?''

Một nốt *LENG KENG* nữa vang lên rồi tắt hẳn, như người gõ cửa chợt nhớ ra mình đi nhầm nhà.

Tĩnh lặng trở lại, sâu hơn cả lúc trước, như có ai đang bịt tai cả thế giới.

Tôi hít một hơi thật sâu, cố giả vờ nhún vai. ``Chắc là gió chợt lùa từ phía kia tường thôi.''

Sakura-san không phản bác. Cả Shino-san cũng không.

Chúng tôi bước tiếp, nhưng mỗi bước đều nặng hơn. Tôi cảm giác có thứ gì vừa gõ nhẹ vào màng thực tại, không đủ mạnh để vỡ, chỉ đủ để in dấu: ``Ta vẫn còn đây.''

Phía sau lưng, chiếc chuông không reo thêm lần nào.

Nhưng tôi biết, những điều không cần âm thanh vẫn có thể thét vào tai ta, nếu ta đủ lặng để nghe.

Hy vọng rằng ngày diễn ra bữa tiệc sẽ diễn ra đầu xuôi đuôi lọt\dots

\begin{center}
  \uline{Hai ngày sau.}
  \pov{Hikawa Suzuki}
\end{center}

Hai ngày trôi qua nhanh hơn tôi tưởng. Để ``báo cáo'' nhanh: suốt 48 giờ ấy chẳng có hiện tượng lạ nào đáng kể. Không tiếng chuông gió thình lình reo, không bóng người lướt qua góc hành lang, cũng chẳng có cơn gió lạ nào cuốn mùi hương xa lạ vào lớp học.

Sáng nay vừa đặt chân vào lớp, tôi đã cảm được một luồng khí khác hẳn mọi ngày. Shino-san ngồi thẳng băng, hai tay siết lấy mép bàn như sợ bị cuốn đi; Sakura-san thì cứ mỗi năm giây lại liếc đồng hồ treo tường, rồi lại nhìn ra cửa sổ, chờ đợi điều gì chỉ mình cô hiểu.

Sáu bạn chuyển trường cũng không giấu được vẻ hồi hộp-- à, năm người thôi, Koishin-san vẫn mặt lạnh tanh, nhưng tôi đánh cược là trong lòng cô ấy đang đếm từng nhịp thở.

Tôi lẽo đẽo tới chỗ ngồi quen, kế bên Onee-sama, thì thấy Onii-sama đang cười đùa với Akane-san và Shino-sam. Anh trông thì thản nhiên, nhưng tôi biết tai anh dựng đứng, theo dõi từng tiếng động trong lớp.

Chẳng bao lâu, tiết Sinh hoạt lớp cuối buổi bắt đầu.

Tsukishita Yuko-sensei bước vào lớp, gõ nhẹ thước lên bàn để ra hiệu trật tự. Cô mỉm cười hiền, đủ để xoa dịu, nhưng đáy mắt sâu thẳm khiến ai nấy thắt lại: chắc chắn sắp có chuyện.

``Hôm nay cô có thông báo quan trọng.''

Vài bạn đã nín thở. Yuko-sensei khẽ gật: ``Cuối buổi nay, lớp 1-C sẽ tự tổ chức một buổi gắn kết.''

Cả lớp chết lặng một nhịp. Nhưng tôi, Onii-sama, Onee-sama và mấy bạn mới đã ngờ ngợ từ trước.

``Ý tưởng này,'' cô tiếp, ``được đề xuất bởi Shirayuki Inari, Natsuno Hanabi, Tokiwa Kana và Amano Kanade.''

``\textbf{Thật hả!?}'' Kaoru-san bật dậy khỏi ghế.

``Có ăn uống không?'' Hotaru-san hỏi sốt sắng, mắt sáng như đèn pha.

Ayame-san chống nạnh: ``Có tiệc thì phải có trò vui chứ!''

Honoka-san suýt reo: ``Trời, nghe đã thấy hào hứng rồi\~{}!''

Tiếng bàn ghế xê dịch, tiếng cười thì thầm hòa vào nhau, biến căn phòng quen thuộc thành một lò xo ấm áp.

Onee-sama cười rất thoải mái, kiểu cười mà tôi chỉ thấy khi chị thật sự thả lỏng.

Onii-sama thì lặng lẽ quan sát, nhưng khóe môi anh hơi cong lên một chút. Rất nhẹ, nhưng tôi nhận ra được.

Tôi siết nhẹ tay trên đùi, mỉm cười.

Và tôi mong rằng, không chỉ cho Shino-san, Sakura-san hay lũ bạn mới, mà cho cả lớp 1-C, đây sẽ là khởi đầu thật sự của một thế giới không còn quay vòng. Bữa tiệc này sẽ trở thành ký ức không thể quên.

\ct

Buổi Sinh hoạt lớp vừa kết thúc, cả lớp vẫn còn rì rầm bàn tán thì Yuko-sensei đã vỗ tay một cái, kéo mọi người quay lại trật tự.

``Vì buổi tiệc sẽ diễn ra chiều nay,'' cô nói, ``nên buổi sáng này lớp ta sẽ chia nhóm để chuẩn bị. Cô sẽ không can thiệp sâu, các em tự bàn bạc và chịu trách nhiệm.''

Câu nói đó như bật công tắc.

Chưa đầy vài phút sau, bảng phân công đã được thống nhất, nhanh đến mức tôi còn hơi ngạc nhiên.

Azuki-san, Yuki-san, Mitsuki-san và Saya-san cùng nhau đảm nhận khâu âm nhạc, còn Sakura-san thì tự nguyện gánh thêm trống nhỏ cho bữa tiệc thêm rộn ràng.

Phần trang trí được giao cho Suzune-san, Inari-san, Hotaru-san và Wata-san, bốn cô gái hẳn sẽ biến lớp học thành một góc trời ấm áp.

Ayame-san, Kaoru-san, Miku-san, Kana-san và Kanade-san lập tức xung phong lập tổ trò chơi, hứa hẹn mang đến tràng cười không dứt.

Shino-san, Onee-sama, tôi và Mari-san tụ lại thành nhóm bếp, lo từ khâu nấu nướng đến thổi lửa.

Cuối cùng, Onii-sama được ``vinh danh'' làm tổng hậu cần, chuyên lo chạy chợ rút ruột để không thiếu cái gì. Nghe xong phân công, anh quay lại nhìn bảng, cứng đờ đúng hai giây.

``\dots Sao lại dính tới tôi?''

``Ai bảo cậu rảnh nhất lớp làm chi\~{}'' Ayame-san đáp ngắn gọn, không thèm ngại.

``Cậu còn là hội viên danh dự của Câu lạc bộ Về nhà mà.'' Kaoru-san bồi thêm một câu nghe rất logic\dots\ nhưng lại chẳng logic chút nào.

Onii-sama thở dài. Nhưng tôi biết, khi đã bị kéo vào rồi thì anh ấy sẽ không bỏ giữa chừng.

Tôi đứng lùi lại một chút, quan sát từng nhóm bắt đầu vào việc.

Ở góc phòng, Azuki-san đang khởi động giọng. Chỉ là vài nốt thử thôi, vậy mà cả lớp đã ồ lên khe khẽ. Giọng cô ấy vang và chắc, không cần phô trương nhưng khiến người ta tự động chú ý.

Saya-san đứng cạnh, nhanh chóng chuyển sang vai trò điều phối. Cô ấy phân chia thứ tự tập, nhắc nhở vị trí đặt nhạc cụ, còn ghi chép lại những gì cần chuẩn bị thêm. Nhìn cách cô ấy làm việc, tôi có cảm giác rất yên tâm.

Trong khu bếp tạm thời của lớp (thực chất là một góc lớp học với những dụng cụ được mượn từ câu lạc bộ Nấu ăn), Onee-sama đã xắn tay áo từ lúc nào. Dao trong tay chị chuyển động dứt khoát, đều và nhanh. Cà rốt, hành, khoai\dots\ tất cả được xử lý gọn gàng chỉ trong nháy mắt. Tôi đứng cạnh rửa rau, liếc sang mà không giấu được chút tự hào lẫn xót xa: ở thế giới cũ, khi cả hai chị em còn chật vật từng bữa, thực phẩm tươi như thế này là thứ xa xỉ chỉ dám mơ tới. Những lần hiếm hoi kiếm được củ cà rốt không dập, chị cũng thái thật mỏng để kéo dài, từng đường dao đều tăm tắp như bây giờ, chỉ khác là ngày ấy chúng tôi ăn chung một bát, chia nhau từng miếng.

Shino-san và Mari-san thì lo phần nêm nếm, thỉnh thoảng trao đổi nhỏ với nhau, giọng nhẹ nhưng rất tập trung.

Ở giữa lớp, tổ trò chơi là nơi ồn ào nhất.

Ayame-san dựng lên một cái bảng lớn, treo đầy chuông nhỏ đủ màu. ``Trò này gọi là đập chuông bingo!'' cô ấy cười toe. ``Trúng là có thưởng liền!''

Tôi nhìn mà toát mồ hôi. ``C-Cái này có hợp với tiệc chào mừng không vậy!?'' tôi buột miệng.

``Yên tâm đi!'' Hội trưởng hội học sinh khoát tay. ``Cũng giống như tổ chức lễ hội đó thôi, yo!''

Bên cạnh, Kana-san và Kanade-san đang tranh luận rất nghiêm túc về giải thưởng.

``Giải nhất phải là thứ gì đó thật hoành tráng!''

``Không được, phải công bằng cho mọi người chứ!''

``Nhưng nếu nhạt quá thì ai hứng thú? Bà có cao kiến gì khác không?''

Cuộc tranh luận nghe chẳng khác gì một hội đồng thu nhỏ.

Còn Onii-sama, tuy mặt cau có như bị ép chợ, vẫn lọ mọ giữa các nhóm, vừa đọc lô-tô vật dụng vừa ra lệnh ``chỗ này thiếu, chỗ kia thừa'', xếp lịch từng khung giờ như thể đang điều hành công trường.

Tôi bỗng nhớ ra: anh ấy từng chia sẻ rằng là một tổng quản tại một công ty giải trí, hóa ra cái máu quản lý nó ngấm vào xương rồi, bỏ không được!

Buổi sáng trôi qua trong nhịp độ nhanh và hơi hỗn loạn, nhưng không hề khó chịu. Ngược lại, có một cảm giác gì đó rất thật, như thể lớp 1-C đang cùng nhau tạo ra một điều chung, ngay tại đây.

Tôi hít một hơi nhẹ, rồi tiếp tục công việc của mình.

Shino-san lặng lẽ đứng bên bếp, khéo léo xếp từng khay bánh mới ra lò. Dù cô ấy không nói, tôi vẫn cảm nhận được niềm hạnh phúc kín đáo qua từng động tác: kéo cho khăn trải bàn thật phẳng, nghiêng khay kiểm tra bánh đã thẳng hàng chưa. Có lẽ với cô ấy, đây chính là giây phút từng tưởng chừng xa vời: được cùng lớp chuẩn bị bữa tiệc thân thuộc, không vương vấn bất cứ định mệnh u ám nào.

Cả lớp 1-C bỗng biến dạng. Bàn ghế xếp dồn hai bên, giấy màu và ruy-băng rủ xuống lơ lửng; mùi ngọt quyện với tiếng cười nói rộn rã. Góc kia, nhạc công thử âm lần cuối; ngoài hành lang, vài bạn hối hả mang đồ trang hoàng. Tất cả nhịp nhàng, rộn ràng, như mỗi người đang thêm một mảnh ghép cho kỷ niệm chung sắp thành hình.

Hai chúng tôi mỗi người bưng một đĩa bánh ngọt, chuẩn bị mang ra bàn phía trước lớp.

Nhưng rồi, bất chợt\dots

Một bóng dáng nhỏ vụt qua hành lang.

Tôi khựng lại theo phản xạ. Trong khoảnh khắc rất ngắn ấy, tôi thấy rõ ràng một đứa trẻ thấp bé, mặc kimono cũ rách, tóc xõa rối, chạy ngang qua tầm mắt chúng tôi. Không có tiếng bước chân. Chỉ là một cái lướt qua, nhanh đến mức nếu không nhìn thẳng, tôi đã nghĩ mình tưởng tượng ra.

``\dots!''

Shino-san khẽ hít mạnh một hơi. Đĩa bánh trên tay cô ấy chao đi, suýt nữa thì trượt khỏi lòng bàn tay.

``Ơ-- Shino-san, cẩn thận!'' Mari-san thảng thốt kêu lên từ phía sau.

Tôi vội vàng đưa tay đỡ, giữ cho đĩa bánh không rơi. Tim tôi đập thình thịch, cổ họng khô lại. Tôi quay đầu nhìn theo hướng hành lang vừa rồi, nhưng nơi đó đã trống không. Chỉ còn lại cánh cửa góc cầu thang khép hờ, im lìm như chưa từng có ai đi qua.

Cô nàng tóc nâu đứng bất động vài giây, mắt vẫn dán vào hành lang vừa vụt qua bóng kimono. Mari-san giật mình hỏi: ``C-Cậu làm sao thế, Shino-san?''

Shino-san khẽ lắc đầu, giọng thì thầm như sợ ai nghe thấy: ``T-Tớ ổn. Và\dots\ cảm ơn cậu.'' Dù vậy, tôi thấy rõ bàn tay cô khẽ siết chặt mép đĩa đến trắng bệch, đầu ngón tay run nhè nhẹ.

Tôi nuốt khan, tự nhủ trong đầu: ``\textit{Không biết có phải là ảo giác không nữa.}''

Chúng tôi vẫn tiến về phía trước, đặt từng đĩa bánh lên bàn như kế hoạch. Không khí lớp học nhanh chóng trở nên sôi động, tiếng cười nói tràn ngập, xua tan đi vẻ yên ắng ban nãy. Không ai khác để lộ dấu hiệu nhận ra điều gì lạ.

Dù bề ngoài đã trở lại nhịp sinh hoạt thường ngày, tôi hiểu, và Shino-san cũng thế, rằng trong lòng chúng tôi vừa lướt qua một gợn sóng nhỏ. Mỏng manh, mơ hồ, song đủ khiến nỗi lo âu len lỏi, như bóng đen mờ nằm sát mép bữa tiệc đang đến gần.

% Buổi chiều nay chắc chắn sẽ rất nhộn nhịp lắm đây.

\ct

Buổi chiều buông xuống bất ngờ, nhanh như thể ai đó vội vàng kéo rèm che mặt trời. Tia nắng cuối ngày loang dần trên hành lang, phủ một lớp vàng úa mỏng tang, như màu của ký ức sắp phai. Chúng tôi tập trung lại hết vào lớp 1-C, tiếng bước chân rộn rã hòa vào nhịp tim đập thình thịch trong lồng ngực.

Giờ đây, lớp học mở ra, đón một dòng người đầy màu sắc: những sợi ruy-băng còn vương mùi giấy mới, bóng bay lơ lửng như mới thoát khỏi bàn tay vội vã của ban trang trí, và mùi thức ăn với đầy đủ những vị cay, ngọt, béo, nồng; tất cả quyện lại thành một bản giao hưởng khiến bao tử tôi réo lên một cách thật thà.

Inari-san đứng trên bục giảng, đuôi cáo (phụ kiện, tôi biết) khẽ rung như chiếc vương miện tự phong. Cô vỗ tay hai cái, đôi mắt cười cong cong, rồi cất giọng trong veo vang khắp gian phòng: ``Mọi người ơi! Bữa tiệc chào đón học sinh mới của lớp 1-C chúng ta\dots\ \textbf{chính thức bắt đầu!}''

Tiếng reo hò bùng lên như pháo tết. Hanabi-san bật chiếc pháo giấy trên tay, và *BỐP!*, những mảnh bóng bay nhỏ xíu tung tóe như hoa gạo rơi. Cô xoay một vòng, váy xòe tung, cười đến mức cả hành lang dường như cũng phải nghiêng ngả theo.

``\textbf{Làm tới đi nàooo\~{}!}'' cô hét vang, giọng vỡ òa trong tiếng vỗ tay rộn rã.

Chưa kịp dứt nhịp, Ayame-san đã lướt từ đâu tới, chuông kim loại trên tay rung lên chan chát. Cô đập bàn trò chơi một cái *keng*, rồi nheo mắt: ``Sau lễ khai mạc, nhớ xẹp qua gian hàng `Đập chuông Bingo' của tôi nha! Có bánh, có kẹo, có cả vé miễn phát ngôn bậy một ngày đóooo!''

``Ayame!'' Kana-san đỏ mặt, giậm chân. ``Chưa tới phần trò chơi mà! Đừng có phá hoại chương trình!''

``Quảng cáo sớm mới giật khách chứ!'' Hội trưởng thè lưỡi, còn cố tình đập thêm một phát chuông nữa cho đã nghiệp.

Tôi bật cười, tiếng cười thoát ra như bong bóng xà phòng. Trong khoảnh khắc ấy, mọi lo toan ban nãy về vòng lặp, về bóng kimono kia, về tiếng chuông gió không gió, như tạm thời bị xua tan bởi một làn khói thơm mỏng. Tôi nhìn quanh: Onee-sama kéo ghế xuống cho Onii-sama ngồi, Shino-san e ấp cầm khay bánh, còn Sakura-san thì len lén chỉnh lại dải ruy-băng đã bị kéo lệch. Tất cả đều hiện hữu, đều ấm, đều thật.

Và tôi biết, dù chỉ là một buổi chiều, dù ngày mai có thể mang đến những điều không ai ngờ, thì khoảnh khắc này, khi tiếng cười còn nồng, khi mùi thức ăn còn nóng, khi trái tim tôi đập đều nhịp với cả lớp, đã là một khởi đầu thật sự.

Bữa tiệc đã thật sự bắt đầu rồi.

\ct

Những chiếc bàn được xếp thành hình chữ U, phủ khăn trắng tinh, trông như một dải lụa vắt ngang lớp học.
Tổ nấu ăn bọn tôi tự thưởng cho mình điểm mười vì màn trình diễn này.

Sandwich trứng cắt tam giác của Shino-san, với viền bánh căng mịn, lớp nhân vàng ươm như những khối rubi nhỏ xíu, đẹp đến mức tôi phải chụp ảnh trước khi dám gắp.

Súp kem bắp mà Onee-sama khuấy, mùi ngọt béo bốc lên nghi ngút, mặt súp phẳng lỳ như tấm gương, in bóng cả dãy đèn neon trên trần.

Gà rán sốt mật ong của Mari-san lớp vỏ vàng ươm giòn tan, phủ lớp sốt mật ong óng ả, sáng bóng như mặt hồ trong nắng chiều.

Và bánh trái cây của tôi, cùng với sự giúp sức của Wata-san: cốt bánh mềm xốp, phủ dâu tây, kiwi, đào thành vòng hoa nhiệt đới, quệt thêm lớp kem tươi bồng bềnh như mây.

``Woa\dots\ đây là tiệc hay bảo tàng ẩm thực vậy?'' Hotaru-san che miệng, mắt sáng như đèn pha.

``Kaigou-san, cậu mở nhà hàng đi, mình xin làm khách quen!'' Wata-san nắm chặt tay, cười tít mắt.

Chị tôi hất mái tóc, nháy mắt: ``Nếu mở thì chỉ phục vụ nữ thôi nhé\~{} Các chàng trai đứng ngoài ngửi mùi thôi!'' Dĩ nhiên là cũng có ngoại lệ, và chắc các bạn cũng biết đó là ai.

Onii-sama đứng từ phía cửa lớp, tay khoanh, mắt đảo qua từng khay như radar: ``Khay nhỏ thế này thì lấy đâu đủ thìa đây? Đổi sang khay lớn, xếp thành hai hàng cho dễ lấy.''

Dù miệng càu nhàu, anh vẫn tự tay sắp lại từng cái muỗng, chỉnh từng góc khăn cho vuông vức.

Tôi mỉm cười, thầm nghĩ: ``\textit{Onii-sama vẫn vậy, lạnh lùng bên ngoài nhưng bên trong là cái máy tính siêu chi tiết.}''

``Này, bánh này nhìn như tác phẩm nghệ thuật ấy!''  Nanase-san xuất hiện bên cạnh, chỉ vào khay bánh trái cây. ``Cắt một miếng thôi cũng thấy có tội với tác giả.''

``Đừng lo, tác giả thích người ăn hết đấy,'' tôi đáp, đẩy khay về phía cô. ``Miễn là cậu ăn hết phần là được.''

Kaoru-san len lén gắp một miếng gà, cắn giòn tan, mắt sáng lên: ``Mari-san, cậu bỏ bí kíp gì trong đây? Vỏ giòn mà thịt vẫn ẩm, ngọt thấm đều!''

Cô bạn tóc đỏ che miệng cười: ``Bí mật của nhà Kaigou, kể ra mất thiêng.''

Shion-san đứng xa xa, ôm cánh tay, nhưng mắt không rời khay sandwich: ``Hừm, nếu ta thêm chút bùa may mắn vào, chắc sẽ ngon hơn\dots\ nhưng thôi, để lần sau.''

Uzuki-san thì đã kịp xếp hàng hai lần: ``Mình vừa ăn thử súp rồi, mịn như lụa, ngọt thanh như hơi thở mùa thu! Kaigou-san, cậu có dạy nấu ăn không? Mình muốn đăng ký học viên đầu tiên!''

Tôi lau tay vào tạp dề, quay sang người đang dựng thẳng thớm ở cửa lớp như cột điện. ``\vie{Onii-sama, vào ăn chung đi!}''

Anh lắc nhẹ đầu, mắt vẫn đảo qua từng khay. ``\vie{Anh sẽ ăn sau, mọi người cứ ăn trước đi.}''

Chưa kịp đáp, Onee-sama đã cất tiếng châm chọc: ``\vie{Không vào bây giờ thì e rằng mấy cái đùi gà sẽ chỉ còn mỗi xương, còn súp thì bay hơi hết sạch đấy nhé, Kuon-kun\~{}}''

Cả lớp bật cười rộn rã; tiếng cười như thác đổ làm cậu ta khựng nửa giây rồi chỉ nhún vai, mặt như viết hẳn dòng chữ ``kệ tôi''.

Không khí lớp học bỗng như được thắp lửa: những lời khen ngợi, tiếng cười giòn tan, tiếng thìa đũa gõ nhịp vào khay inox tạo thành bản giao hưởng rộn rã. Tôi nhìn quanh, thấy bức tranh ấm áp đang dần hoàn tất: mùi bơ nóng quyện mật ong, sắc đỏ của cà chua bi, xanh của lá húng, vàng của đùi gà chiên, tất cả hòa vào những nụ cười chưa từng được nở trọn vẹn như thế. Đây sẽ là ký ức đầu tiên (và cũng có thể là ký ức cuối cùng) của lớp 1-C nếu thế giới lại quay vòng.

Ánh đèn neon trên trần khẽ lùi sáng, nhường chỗ cho hai dây đèn xanh dương được Kanade-san quấn quanh cửa sổ, tạo thành một khung sân khấu thu nhỏ. Azuki-san bước ra, tay cầm micro không dây, mắt khép hờ như đang nghe tiếng gió thở trong huyền thoại. Yuki-san đứng cạnh, đàn ghi-ta ôm ngang bụng, móng tay lướt nhẹ lên dây để thử âm. Một tiếng *tinh* vang lên, căng phòng như mặt hồ được gõ nhẹ bằng ngón tay thủy tinh.

Giai điệu mở đầu. Một điệu ballad pha chút bossa tràn ra. Cô nàng cầm mic cất tiếng: trong veo nhưng không mỏng manh, mỗi chữ như hạt mưa đầu mùa rơi xuống lá sen. Cùng lúc đó, Yuki-san đẩy thanh lên quãng cao, giọng anh mềm như tuyết tan, nhưng ẩn sau đó là một lực kéo vô hình khiến ngực người nghe thắt lại. Cả lớp im phăng phắc, ngay cả Ayame-san cũng quên đập chuông quảng cáo.

Tôi thấy mũi mình châm chích. Không dám chớp mắt, tôi nghiêng đầu, và thấy Shino-san đã lặng lẽ đứng sát vai tôi từ lúc nào. Bàn tay cô siết nhẹ tay tôi, đầu ngón tay lạnh toát.

Đúng lúc ca khúc leo tới đoạn cao trào, đèn trần trên phòng bất chợt nhấp nháy một cái, chỉ như có người vừa vẫy màn đêm qua khung cửa. Một tiếng *ROẸT* vô hình đâm thẳng vào võng mạc. Tất cả thở gấp một nhịp. Rồi đèn sáng lại, như chưa từng biến động.

``Cầu dao giật thôi mà!'' Kana-san liền cười khẽ, nhưng nụ cười cô cứng đờ trước khi kịp hoàn tất.

``Đúng rồi, đúng rồi\dots'' Inari-san vội đáp lời, đuôi cáo phụ kiện rung lên bất an.

Tôi quay sang Shino-san. Mắt cô mở to, đồng tử co lại. Cô không nhìn tôi, mà nhìn *xuyên qua* tôi, vào khoảng không sau lưng. Tôi bám theo tầm mắt cô.

Những đốm sáng.

Chúng lơ lửng giữa không trung, to bằng hạt đậu, phát ra thứ ánh sáng mờ như bụi phản quang. Nhưng không phải bụi thông thường, mà chúng có hình dạng hẳn hoi: đầu tròn, tay ngắn, chân khấp khểnh. Những đứa trẻ thu nhỏ, chỉ bằng ngón tay cái, đang nghiêng đầu, khẽ đung đưa theo nhịp của bản ballad. Không tiếng thở, không cánh, chỉ lấp lánh như đom đóm mùa thu, nhưng lại không có khuôn mặt người.

Shino-san siết tay tôi chặt hơn. ``Tớ\dots\ thấy rồi,'' cô thì thầm, giọng khô khốc.

Tôi nuốt khan. Cổ họng như dính bột. Không phải ảo giác, cũng không phải đèn led bị lỗi. Những vị khách vô hình đang \emph{nghe nhạc} cùng chúng tôi, bình thản như thể đây là buổi hòa nhạc dành riêng cho họ.

Cặp đôi trên sân khấu vẫn hát, mắt nhắm nửa, môi mỉm cười. Họ không biết những gì đã xảy ra, hoặc có lẽ,  đã chọn vờ như không biết. Tôi đứng chết lặng, chỉ kịp thì thầm một lời cầu trong đầu:

\begin{center}
  ``\textit{Xin hãy để bữa tiệc này trọn vẹn.
    Xin đừng bắt chúng tôi quay lại từ đầu\dots\ thêm một lần nữa.}''
\end{center}

\ct

Tiếng nhạc vừa tắt hẳn, còn chưa kịp tan trong không trung thì Ayame-san đã nhảy phốc lên bục giảng, tay quăng micro đồ chơi lấp lánh ánh bạc như món đồ chơi mua vội ở hội chợ đêm. Cô xoay một vòng, váy đồng phục xòe tung, rồi hét vang bằng giọng MC đầy lầy lội: ``\textbf{Chào mừng quý vị đến với màn TRÒ! CHƠI! đỉnh cao của lịch sử lớp 1-C! Ai không hét to, coi như tự nguyện vào danh sách đen của tôi đó nhaaaa~!}''

Toàn bộ lớp vừa cười vừa rụt cổ. Onii-sama đứng cạnh tôi, khoanh tay, khẽ thở dài đủ để tôi nghe thấy tiếng cười chát chúa trong cổ họng anh: ``Cứ như thể chúng ta đang ở lễ hội mùa hè ấy nhỉ.''

Hội trưởng hội học sinh lập tức quay ngoắt 180 độ, nheo mắt, chỉ thẳng mũi anh: ``Kuon-kun, cậu còn thở dài thêm một cái nữa thì tôi cho cậu làm đội trưởng đội thua cuộc đó! Đừng có khóc ròng khi bị Shino-san vượt mặt nhé!''

Tôi bật cười sặc sụa, nhưng trong đầu vẫn lởn vởn những đốm sáng tí hon lúc nãy: những ``vị khách'' không mời mà đến đang lơ lửng giữa trần như đom đóm lạ. Tôi lén siết chặt tay áo, tự nhủ với bản thân: ``\textit{Đừng để chúng phá bữa tiệc\dots}''

Chưa đầy một phút, cả lớp đã được chia thành hai đội dàn hàng ngang.

Đội Onee-sama (tên gọi không chính thức: ``Hội những người nấu ăn quá giỏi'') gồm: Onee-sama, Yae, Koishin và Wata, bốn cái tên vừa nhìn đã ngửi được một mùi đồ ăn thơm phức bay ra.

Đội Onii-sama (biệt danh tạm thời: ``Hội anti-lãng phí thời gian'') gồm: Onii-sama, Shino, Nanase và Kana, bốn cái tên vừa nhìn đã thấy họ quyết tâm không thừa một giây.

Còn tôi? Tôi bị kẹt ở giữa, vừa được Ayame-san xỏ ngay cái bảng ``TRỌNG TÀI DỰ BỊ'' to tướng, vừa bị cô ấy thúc vai: ``Suzuki-chan, nhiệm vụ cao cả của cậu là canh me xem đội nào ăn gian thì báo ngay cho tôi! Từ đó, tôi sẽ có biện pháp xử lý mạnh tay nha\~{}!''

Tôi mở to mắt, rồi bật cười gật đầu. Giữa bao điều kỳ lạ chưa có lời giải và vòng quay thời gian chưa biết điểm dừng, khoảnh khắc cả lớp cùng reo hò vì một trò chơi nhỏ bé cũng trở nên quý giá vô cùng.

Trò chơi mang tên ``Bắt thẻ bài sắc màu'' được mở màn. Luật chơi thật dễ: sàn lớp rải đầy những tấm thẻ đủ màu; đèn neon bật màu gì, người chơi phải vội vàng nhặt đúng màu ấy. Kẻ chậm tay hay sai màu sẽ lập tức bị loại.

Ngay khi dòng điện màu đỏ bùng lên trên thanh neon, cả lớp như bị ai giật dây. Onee-sama phóng ra trước tiên, tay phải chém không khí, bàn tay trái chụp lấy tấm thẻ đỏ như hái một bông hoa chớp nở. Koishin-san theo sát, khịt mũi như mèo ngửi mồi, vớ được thẻ xong lập tức lùi về sau hàng, miệng lẩm nhẩm ``Tôi không cố đâu, tôi chỉ\dots tình cờ thôi'' để bảo toạn danh dự ``lạnh lùng'' của mình.

Đối diện, Onii-sama không thèm nhúc nhích. Anh đứng yên vj, mắt đo khoảng cách, chờ thẻ rơi đúng tầm rồi khẽ cúi nhặt như nhặt lá vàng trên mặt hồ. Nanase-san reo hò vỗ tay ``Đúng phong cách Kuon-san!'', còn Kana-san mải ngắm giải thưởng ``vé miễn phát ngôn bậy'' nên suýt nhặt nhầm thẻ tím, may có Shino-san kéo lại kịp. Tiếng cười, tiếng la hò, tiếng dép lê cọ xát sàn gỗ tạo thành một bản giao hưởng đầy nhựa sống. Tỉ số 1-1, căng như dây đàn ngay từ lượt thi đấu đầu tiên.

Chưa kịp lấy hơi, đèn neon bật xanh dương. Shin-san lao đi như làn khói, tóc cô vẩy ra sau, chỉ để lại hương hoa nhài thoang thoảng. Cô gắp thẻ xong dừng lại, khẽ thở, mắt ánh lên niềm nắng nhỏ, có lẽ lần đầu cô thấy mình ``chạy'' vì điều gì đó không phải nỗi sợ.

Ayame-san giơ micro nhựa, hét vỡ cổ họng: ``WAO! Shino nhà ta đã vào form! Còn ai dám thách không?'' Onee-sama cười khẽ, gật gù: ``Được đấy, em gái tôi cũng không kém.''

Thẻ bài lúc này như lá thu rơi, lách cách bay khắp lớp. Có đứa trượt ngã, có đứa ôm bụng cười, có đứa la ``\textbf{Màu vàng! Vàng đâu rồi!}'' rồi vớ vội tấm cuối cùng. Tôi đứng giữa, bỗng thấy lòng mình nhẹ bẫng. Những đốm sáng kỳ lạ kia vẫn lơ lửng trên cao, nhưng chúng chỉ nhấp nháy theo nhịp nhạc, không dám xuống. Có lẽ chúng cũng đang cảm thấy vui?

Trong tích tắc ấy, tôi như bị búng ra khỏi dòng thời gian: dù ngày mai đồng hồ có quay ngược, dù bữa tiệc này bị xóa sạch như bụi trên bảng, thì hiện tại—tiếng cười nức nở, mùi bánh nóng hổi, ánh đèn neon loang loáng và cả những bước chân vội vã—đã khắc vào tôi một vết ấn không gì xóa được. Tôi đang sống trong giấc mơ mà trước đây chỉ dám thấy qua khe cửa hé. Và tôi thề sẽ giữ lấy, dù chỉ bằng một tấm thẻ màu xanh dương nhỏ bé đang nhúng mồ hôi trong tay.

Nhưng rồi\dots\ mọi thứ như bị ai đó bấm nút tua nhanh.

Đèn neon lại thay đổi, màu tím sậm, màu của hoàng hôn chết.

Một tấm thẻ tím rơi đúng trước mặt tôi, lượn một vòng như lá thu cuối cùng.

Tôi cúi xuống định nhặt lên\dots\ thì một bàn tay nhỏ xíu, lạnh buốt như viên đá vừa mới bóc từ hầm băng, chạm vào trước.

Tôi khựng lại, hơi thở tắc nghẽn trong cuống họng.

Một đứa trẻ đứng trước mặt tôi. Kimono rách sờn, đường chỉ đỏ chạy loang lổ như vết máu khô. Không có khuôn mặt, chỉ là một mặt phẳng trắng bệch, nhưng tôi lại cảm giác rõ ràng nó đang nhìn mình bằng một đôi mắt vô hình.

Nó cúi xuống, nhặt thẻ tím lên. Và rồi chậm rãi, như người ta đặt một bông hoa lên mộ, nó đặt vào tay tôi.

Không ai khác phản ứng. Không ai nhìn thấy.

Chỉ tôi và Shino-san; tôi cảm thấy cô siết cánh tay tôi, móng tay như muốn đâm thủng da.

Tôi nghe giọng mình bật ra, run rẩy như lá cuối cùng trên cành:

``\dots Cảm ơn\dots?``''

Trong thoáng chốc, đứa trẻ tan biến, chỉ còn lại hơi lạnh quanh cổ tay tôi như vệt bụi sáng bị gió cuốn vào đêm.

Mọi thứ trở lại như chưa từng xảy ra: tiếng cười, tiếng dép lê cọ sàn, mùi bánh vẫn ngọt lịm.

Tôi vẫn đứng đó, tay cầm tấm thẻ tím, tim đập mạnh đến mức tưởng như cả lớp đều nghe thấy\dots\ hoặc đang nghe thấy.

``Suzuki-chan, cậu ổn chứ?`` Hotaru-san chạm nhẹ vai tôi, ánh mắt lo lắng như mẹ tôi ngày xưa.

``Cậu trông có vẻ hơi lơ đễnh quá đó na,`` Mitsuki-san ghé vào thì thào, giọng cô như tiếng lá xào xạc.

Yae-san nghiêng đầu: ``Nếu cậu mệt thì ngồi nghỉ cũng được. Đừng cố quá.``

Shiina-san chỉ lẳng lặng quan sát, mắt cô ấy như máy quét đang cố nối ghép điểm gợn sóng giữa hai thế giới.

Tôi lắc đầu, cố mỉm cười, dù môi tôi cứng đơ như đã đóng băng: ``M-Mình không sao\dots\ Chắc chỉ hơi phân tâm thôi.``

Nhưng khi tôi nhìn lại, tôi bắt gặp ánh mắt Onii-sama, một tia sáng cảnh giác im lặng, như con mèo đêm vừa nghe tiếng chuột trên xà.

Onee-sama cũng nhìn tôi, ánh mắt sắc bén như muốn hỏi ``em có thấy nó không?`` mà không cần mở miệng. Chị khẽ gật, như thể đã đọc được câu trả lời trong mắt tôi.

Và ở góc xa hơn, Yuka-san đứng chống cằm, vẻ mặt trầm ngâm như nhà toán học vừa phát hiện ra biến số mới.

Ayame-san thì\dots\ khác hẳn mọi người: đôi mắt hội trưởng ấy dõi theo tôi, một nụ cười mơ hồ khó đoán: không phải cười vui, cũng không phải cười nhạo; nó như bức màn che của ảo thuật gia, phía sau có thể là con át chủ bài, cũng có thể là hố đen vô đáy.

Xung quanh, tiếng reo hò lại nổi lên những đợt sóng, dải đèn neon trên cửa sổ nhấp nháy nhịp nhàng như đang đếm ngược. Hội trưởng tiếp tục tung hứng trò chơi bằng những câu chọc quê vừa lầy vừa duyên.

Nhưng linh hồn tôi đã trượt ra khỏi thân xác một khoảng xa, lơ lửng giữa hơi ấm của bữa tiệc và hơi lạnh từ hành lang vắng.

Đứa trẻ mặc bộ kimono rách ấy\dots\ không mang theo mùi thù hận. Trong khoảnh khắc ngón tay lạnh buốt đặt tấm thẻ tím vào lòng bàn tay tôi, tôi cảm nhận được một niềm khát khao đơn thuần: muốn được tham gia, muốn được chơi cùng. Như thể bao năm thái bị giam giữ trong ký ức tối tăm khiến chúng quên cách cười, chỉ còn biết lặng lẽ theo dõi những bữa tiệc của người sống.

Một luồng điện chạy dọc sống lưng tôi, nửa ấm như ánh nến đầu tiên của buổi chiều, nửa lạnh như gió lạnh từ cầu thang cổ. Trái tim tôi đập loạn nhịp, không còn phân biệt nổi đâu là hiện thực, đâu là hư ảo.

Và tôi chợt hiểu rằng: dù bữa tiệc vẫn rộn ràng, dù tiếng cười vẫn đầy ắp, thì bánh xe vòng lặp đã lăn bánh. Điều gì đó đã bắt đầu, không thể đảo ngược; và trong vũ điệu neon chớp sáng ấy, chúng tôi đang đứng giữa ranh giới của ký ức và hiện tại, của sự sống và những linh hồn chưa kịp lớn.

\ct

Tôi trượt người xuống chiếc ghế cuối dãy, ly nước trái cây đã mất hết đá, đọng lại vài giọt sương lạnh trên mu bàn tay. Không gian như một bể cá khổng lồ: tiếng cười nổi trội, mùi bánh nóng hổi, mùi kem sữa quyện vào nhau, mỗi thứ đều rung lên một nhịp riêng. Ở góc phòng, ba bóng hình còn lại - gồm Koishin-san, Shion-san và Aku-san - đang tụ thành một nhóm nhỏ, tách biệt khỏi dòng chảy chung, như thể họ đang tổ chức một buổi trà riêng giữa lễ hội.

Cô bạn tóc vàng khoanh tay trước ngực, má ửng hồng như người vừa chạy marathon xong: ``Tôi\dots\ không hề mê mấy buổi đông người thế này đâu nhé\dots\ Nhưng mà\dots\ cũng không tệ, cho lần đầu tiên.''

Cô ấy nói xong liền quay mặt đi, như sợ ai bắt gặp nụ cười đang trực trào. Tôi bật cười khẽ. Độ tsundere của cô ấy có lẽ đã thành một dải ngân hà đỏ rực luôn rồi.

Shion-san xoay chiếc que phép giả bằng nhựa lấp lánh, mắt sáng như đèn pha: ``Hế hế\dots\ Bữa tiệc này chính là điểm khởi đầu của Great Sorceress Shion! Từ hôm nay, ta sẽ gieo lời nguyền may mắn xuống lớp 1-C! Phù hộ cho mọi màn thi\dots\ và cả những bài kiểm tra bất ngờ!''

Cô vung que phép, tung tóe những hạt giấy bóng màu bạc, rơi lả tả như tuyết mùa đông. Mấy bạn đứng gần trố mắt, rồi cười ồ lên, coi đó là màn ảo thuật đêm hội.

Aku-san nép sát vào mép bàn, hai bàn tay bám chặt viền áo, ngón cái rung rung: ``Ư-ừm\dots\ tớ vui vì\dots\ không bị ai ghét\dots\ vì trước đây, ở nơi cũ\dots\ tớ thường bị bỏ lại một mình.''

Giọng cô nhỏ xíu, như hạt đường tan chậm trong nước lạnh, nhưng lại đủ khiến không khí xung quanh trở nên dịu hơn một chút.

Tôi mỉm cười, khẽ nghiêng người: ``Không chỉ vui, mọi người đang rất cần ba cậu đấy. Lớp 1-C giờ mới thật sự\dots\ đủ màu.''

\textit{Câu nói ấy\dots\ sao mà quen thế.} Tôi nhớ lại bóng dáng Shino-san khi chạm mặt lần đầu tiên: hồi đó cô nàng tỏ vẻ rất tự ti, lầm lũi ngồi ở một góc, đến mức không dám bắt chuyện với bất kỳ ai, cho tới khi tôi (hoặc đúng hơn, bản ngã Reiko tới).

Và cả\dots\ cả tôi nữa, hồi ở thế giới cũ, khi hai chị em tôi từng sống trong cảnh bị người nhà khinh thường và miệt thị, tới mức không ai thèm ngó ngàng tới chúng tôi.

\textit{Giờ thì cả ba đứa ``bị bỏ lại'' lại đủ màu trong cùng một bức tranh.} Tôi siết nhẹ vạt áo, thầm cảm ơn vòng lặp, dù nó định mang đi bao nhiêu, cũng đã cho ta một buổi chiều được\dots\ không còn trong suốt.

Ngoài khung kính lớp học, ánh hoàng hôn như thể ai vừa đổ tràn lớp mật ong loãng xuống nền trời, lan từng đợt sóng cam nhạt bám vào khung cửa sổ sơn trắng, lên cả bức tường vôi và cả những mái tóc đang rối nhẹ vì một buổi chiều mệt nhoài. Không gian như được quét qua một bộ lọc phim cũ, rung lên từng nhịp chậm rãi, khiến tất cả chúng tôi, dù đang đứng thật sự trước mặt nhau, cũng mang theo hơi ấm phai của ký ức, như thể đây là thước phim đã quay đi quay lại bao lần vẫn không chịu mờ.

Azuki-san duỗi thẳng lưng, ngón tay đan sau gáy, mắt nửa nhắm: ``Tuyệt vời quá đi\~{} Đợt sau mình có thể xin phép tổ chức luôn một đêm nhạc acoustic riêng, đèn dây vàng, bánh quy mật ong, cả lớp nằm dài trên chiếu tatami nghe Yuki-san gảy đàn nhé!''

Wata-san khẽ xoay ngón chân trên sàn gỗ, hai tay chắp sau lưng như mèo vừa giấu móng vuốt: ``Khay bánh bọn mình làm hết veo rồi sao? Vậy là\dots\ mọi người thật sự thích chứ không chỉ khách sáo, phải vậy không?''

Koishin-san phồng má, quay nửa thân, mái tóc tung bay che khuất đôi mắt đang lấp lánh: ``Tôi đâu có khen gì đâu, chỉ\dots\ thấy ăn cũng được, không đến nỗi bỏ thừa.''

Shion-chan chớp chớp hàng mi dày, vẫy que phép nhựa lấp lánh: ``Cảnh báo người yếu tim nha, bánh này ngon quá mức cho phép an toàn, ăn xong có thể bị nghiện đấy!''

Aku che miệng, tiếng cười nhỏ như hạt bông gòn rơi vào nước, mắt cúi xuống nhưng khóe môi cong tít: ``Thôi nào, Shion, bà làm mọi người phát sợ đấy!''

Bên ngoài hành lang dài, Yuko-sensei dựa vào khung cửa, tấm danh sách điểm danh gập làm ba, kẹp dưới cánh tay. Ánh đèn vàng hắt lên gương mặt cô, phủ một lớp bụi thời gian mỏng: ``Các em hôm nay làm tôi nhớ lại buổi liên hoan cuối năm thời tôi còn là học sinh. Tinh thần của lớp 1-C\dots\ khiến cô thật sự tự hào.''

Ayame-san nhảy phóc ra sau lưng cô, hai tay chụm trước ngực như diễn viên đang tạo dáng: ``Và phần lớn là nhờ Hội trưởng Ayame-sama dẫn dắt đúng không ạ!?''

Cô giáo thở dài, vai gồng lên rồi xẹp xuống, nhưng khóe môi vẫn kéo cong, không chịu nổi: ``Ừ, ừ\dots\ một phần cũng là nhờ em đó. Nhưng mà lần sau cô muốn em nghiêm túc hơn.''

``Vâng, vâng\~{}''

Từng nhóm học sinh bắt đầu thu dọn khay bánh, gấp khăn trải bàn, tiếng ghế kéo ken két. Dải đèn huỳnh quang hành lang hắt xuống những vệt sáng dài, kéo bóng ai đó thành hai, rồi ba lần người. Hoàng hôn cứ thế trôi, mật ong cạn dần, để lại một lớp màng vàng mỏng dính vào ký ức ai cần giữ, ai cần quên.

\ct

Chỉ còn ba chúng tôi: tôi - người vừa trải qua cơn rùng mình lạnh giá - Onii-sama với bờ vai vẫn cứng đờ như tấm khiên tàng hình, và Onee-sama, mái tóc dài buông xõa, phủ lên lưng áo một lớp ánh nắng cuối ngày mỏng tang. Không khí trong lớp đặc quánh mùi bánh ngọt đã nguội, tiếng cười nãy giờ như bị hút sạch qua khe cửa, chỉ để lại những chiếc ghế lệch nhau, vài tấm khăn trải bàn nhăn nheo, và dải đèn xanh dương trên cửa sổ đang nhấp nháy yếu ớt, như mạch máu cuối cùng của buổi tiệc vừa tắt thở.

``\vie{Thành công ngoài mong đợi,}'' Onee-sama cất giọng, mắt liếc qua đống khay trống không. ``\vie{Mọi người ăn sạch sẽ, khen liên tục, còn xin công thức. Lần đầu tiên tôi thấy Koishin-san tự nguyện cười.}''

``\vie{Cậu ấy nói `không tệ' ba lần,}'' tôi bật cười, rung rung tấm bảng trọng tài vẫn còn trên cổ. ``\vie{Thế mà còn giấu nữa. Nghe với cậu ấy như bản tuyên ngôn rồi.}''

Onii-sama khẽ hừ, nhưng khóe môi kéo lệch: ``\vie{Azuki-san gửi mail khen ngay trong lúc hát. Bảo `chất lượng cao như khách sạn năm sao'. Cậu ta còn thêm bản nhạc nền nữa.}''

``\vie{Còn em,}'' tôi vòng tay ôm ngực, ``\vie{nghe thấy ba nhóm khác xin phép mượn bếp chiều mai. Bữa tiệc lan truyền rồi này.}''

Onee-sama nháy mắt: ``\vie{Có thể mở tiệc đường dài đến cuối tuần luôn. Nghe đồn cả lớp đang bàn đổi tên thành `Câu lạc bộ ẩm thực 1-C' luôn kìa.}''

``\vie{Đừng,}'' Onii-sama lắc đầu, nhưng giọng nhẹ hẳn. ``\vie{Tôi nghĩ mọi thứ nên dừng lại đây thôi. Đi xa quá kẻo cả trường này biến thành cái hội chợ mini mất.}''

Cả ba chúng tôi cùng phá lên cười, tiếng vang đập vào bức tường trống, xua tan mùi nguội của bánh, của hoàng hôn, của nỗi lo vừa lẩn khuất.

Tôi khom lưng, định nhặt cây kéo rơi giữa bãi chiến trường giấy bóng màu; nhưng hơi lạnh từ phía sau gáy lan xuống sống lưng khiến tôi đông cứng nửa chừng. Không phải gió đêm lọt qua khung kính, mà đó là thứ lạnh của hư vô vừa xé một khe hở trong không gian.

Ở khung cửa lớp, thứ ánh hoàng hôn cam đỏ không hề chạm đất; nó loang qua một bóng hình nhỏ xíu, khiến viền áo kimono rách tung tóe như bị bật ra từ trang giấy cũ. Đứa trẻ không bước tiếp, cũng chẳng lùi, mà chỉ lơ lửng, đôi tay bé xíu giơ ra, ngón tay thon thả như que diêm khuya khoát một khoảng trống về phía tôi.

Khoảnh khắc ấy, tôi nhìn thấy từng đường chỉ đỏ đứt lìa trên thân áo, từng mảnh vải phai màu in hình hoa sakura đã bạc màu; và đôi mắt - nếu có thể gọi là mắt - chỉ là hai hố tròn mờ đục, không tròng, không lòng, nhưng chứa đựng một lời van xin câm lặng: ``\vie{Đừng quay lưng.}''

Tôi nín thở. Trái tim đập mạnh đến mức tôi nghe rõ từng nhịp như tiếng trống thúc vang trong hộp sọ. Một giây kéo dài như cả thế kỷ. Tôi chớp mắt. Bóng hình vỡ tan thành từng hạt bụi sáng, hòa vào làn nắng cuối chiều, không để lại dấu vết, không một âm thanh đứt gãy.

Tiếng thở dài của Onii-sama vang sau lưng tôi, khẽ nhưng đủ đứt gãy mặt nạ bình tĩnh: ``\vie{Suzuki\dots\ em lại thấy nó sao?}'' Giọng anh không ngạc nhiên, chỉ mỏi mệt vì đã đoán trước được câu trả lời. Tôi gật, cổ họng khô khốc.

Onee-sama khép tay lên vai tôi, hơi ấm tỏa ra từ lòng bàn tay như muốn xua tan đốm lạnh vừa bám trên da: ``\vie{Chúng không gây hại, phải không?}'' Câu hỏi nhẹ như lông hồng, nhưng ẩn sau đó là cả một rừng câu hỏi khác: ``\vie{Nếu nếu không thì sao? Nếu lần sau không chỉ là đưa tay?}''

Tôi cắn nhẹ môi, mặn chát vị máu nhỏ: ``\vie{Không\dots\ em cảm nhận được. Chúng đang cố nói với chúng ta, hoặc với riêng em. Như thể\dots\ chúng bị mắc kẹt ở khe hở nào đó giữa hai thế giới, và chỉ có chúng ta mở được lối thoát.}''

Onii-sama quay mặt về phía cửa sổ, ánh mắt anh đụng chạm tới đường chân trời đang chuyển màu chết chóc của hoàng hôn. Giọng anh trầm xuống, từng chữ như rơi xuống từ tầng mây thấp: ``\vie{Thế giới này vẫn còn thiếu mảnh. Câu chuyện lớp 1-C chưa khép lại trọn vẹn. Những linh hồn kia\dots\ có lẽ là đoạn ký ức bị quăng quật, là những `lỗi' mà vòng lặp chưa sửa xong. Giải thoát chúng, cũng là đẩy dòng thời gian ra khỏi vòng lặp vô tận.}''

Gió len qua khe cửa, rèm vải rung lên như cánh tay vô hình vẫy gọi. Tôi chợt thấy rõ: bóng đứa trẻ không chỉ đơn thuần xuất hiện; nó là mảnh ghép thiếu hụt mà mỗi lần quay vòng chúng tôi đều vô tình đạp nát dưới chân. Nếu lần này chúng tôi bỏ qua, có thể ngày mai mặt trời sẽ lại mọc trên một buổi sáng hoàn toàn mới, nhưng cũng hoàn toàn cũ.

Tôi nắm chặt tay, móng ghim sâu vào lòng bàn tay: ``\vie{Vậy thì\dots{} chúng ta sẽ không chạy trốn nữa. Dù phía trước là hành lang tối om của hiện tại hay lớp học tràn ngánh quá khứ, chúng ta sẽ đối mặt để kết thúc đúng nghĩa, để những linh hồn lạc lối được về nhà.}''

Onii-sama khẽ gật, khóe môi anh nhếch lên một đường cong mong manh, đủ để tôi biết, anh đã sẵn sàng. Onee-sama siết vai tôi thêm lần nữa, hơi thở cô ấm áp phả vào tóc tôi: ``\vie{Cả ba chúng ta\dots{} cùng nhau.}''

Ngoài kia, hoàng hôn đỏ ối như than than đang chảy dần trên bầu trời. Ba đứa tôi bước ra khỏi lớp, không ai nóng vội cả. Lần này, chừng nào chúng tôi không buông bỏ để tìm ra sự thật, thì vòng lặp sẽ không bắt chúng tôi quay lại nữa.

\end{document}